\documentclass[11pt,a4paper]{article}
\usepackage[utf8]{inputenc}
\usepackage[T1]{fontenc}
\usepackage[spanish]{babel}
\usepackage{hyperref}
\usepackage{geometry}
\geometry{margin=2.5cm}
\usepackage{parskip}

\title{339 si Louise hubiera tenido alas, la película dura quince minutos}
\author{Alejandro Valenzuela Rivera\\\small Universidad Autónoma de Nuevo León}
\date{\today}
\begin{document}
\maketitle

\begin{quote}
\textbf{Sesgo declarado.} Este trabajo se sostiene desde el sesgo declarado de su autor, y no finge neutralidad. Lo que afirma se puede tumbar, y ese es el punto: el falsificador está escrito, no escondido.
\end{quote}

\section*{Resumen}
Quince minutos dura una historia donde el otro ya tiene el canal: el que no lo tiene es el que necesita el mecanismo.

\section{Por qué se sostiene}
El drama no es la inteligencia del otro: es que no comparte el canal para decir lo que sabe.

\section{De dónde viene}
La autoridad: 'hubiera durado 15 minutos si hubieran tenido ALAS'.

\section{El caso}
Arrival: Louise y los heptápodos, y una lengua que se habita.

\section{Método}
\textquestiondown{}Se le dio el canal, o solo se le habló más fuerte?

\section{Límites, y qué lo tumba}
Un caso donde el otro alcance con el canal que ya trae.

\section{Lo que queda abierto}
\textquestiondown{}El otro tiene el canal, o hay que construirlo?

\section*{Cómo se leyó esta tesis}
Esto se escribió en el repositorio man-made-empathy, del que este candado es una pieza: candados/06 del caso/339 si Louise hubiera tenido alas, la película dura quince minutos.md. La máquina que lo compuso lo hizo desde este mismo archivo, y no agregó nada que no estuviera escrito. Si algo de esto es falso, se dice en el mismo sitio, y se corrige con fecha.

\end{document}
